\documentclass[10pt,a4paper]{article}
\usepackage[utf8]{inputenc}
\usepackage[T1]{fontenc}
\usepackage[french]{babel}
\usepackage{geometry}
\usepackage{amsmath,amssymb,array,longtable,tabularx,multicol,fancyhdr,graphicx,tikz}
\geometry{top=1.05cm,bottom=1.15cm,left=1.15cm,right=1.15cm}
\setlength{\parindent}{0pt}
\setlength{\parskip}{2pt}
\setlength{\emergencystretch}{2em}
\begin{document}
\begin{center}
{\large\bfseries COURS DE MATHÉMATIQUES - Produits remarquables — découverte et application}
\end{center}

\begin{center}\textbf{FICHE PÉDAGOGIQUE}\end{center}
\hrule\vspace{2pt}\hrule\smallskip
\begin{center}\textbf{I. ÉLÉMENTS D'IDENTIFICATION}\end{center}
\small\begin{tabularx}{\linewidth}{@{}X X@{}}
\textbf{Établissement :} Établissement de test & \textbf{Année scolaire :}  \\
\textbf{Discipline :} Mathématiques & \textbf{Date :}  \\
\textbf{Classe :} 4e & \textbf{Effectif :}  \\
\textbf{Nom du professeur :}  & \textbf{Nombre de groupes :}  \\
\textbf{SA :} SA1 & \textbf{Titre SA :} Configurations du plan \\
\textbf{Durée curriculaire SA :}  & \textbf{Séquence :} Séquence 8 \\
\textbf{Titre séquence :} Calculs sur les expressions algébriques & \textbf{Durée de la séance :} 55 min \\
\textbf{Numéro de séance :}  & \textbf{Titre du cours :} Produits remarquables — découverte et application \\
\end{tabularx}
\begin{center}\textbf{II. ÉLÉMENTS DE PLANIFICATION}\end{center}
\textbf{1. Contenus de formation :} Factorisation ; produits remarquables ; développement et réduction.\par
\textbf{2.a Compétences disciplinaires :} Résoudre un problème ou une situation-problème en utilisant les concepts et procédures du langage et du raisonnement mathématique.
Appréhender les mathématiques dans ses aspects numériques par l’appropriation des outils, techniques et procédés conventionnels ainsi que par le traitement des données.
Appréhender les mathématiques dans ses aspects géométriques par l’appropriation d’outils et de démarches propres à la géométrie.
Appréhender le calcul sur les expressions algébriques par l’appropriation des outils, techniques et procédés conventionnels de calcul ainsi que le traitement des données.\par
\textbf{2.b Compétence transdisciplinaire :} \par
\textbf{2.c Compétences transversales :} Résoudre un problème ou une situation-problème en utilisant les outils et propriétés relatifs au calcul sur les expressions algébriques.\par
\textbf{3. Connaissances et techniques :} Reconnaître une somme et un produit ; factoriser par un facteur commun ; énoncer et utiliser les produits remarquables ; réduire une somme ; développer un produit ; calculer une valeur numérique.\par
\textbf{4. Stratégie objet d'apprentissage :} \par
\textbf{5. Durée :} \par
\textbf{6. Stratégies d'enseignement/apprentissage :} \par
\textbf{7. Matériels apprenants :} \par
\textbf{8. Matériels enseignant :} \par
\textbf{Instructions du guide mises en œuvre :}\begin{enumerate}
\item Énoncer les propriétés des produits remarquables : $(a+b)^2=a^2+2ab+b^2$, $(a-b)^2=a^2-2ab+b^2$ et $(a+b)(a-b)=a^2-b^2$.
\item Utiliser les produits remarquables pour factoriser, développer et effectuer des calculs numériques.
\item Développer un produit.
\end{enumerate}
\begin{center}\textbf{III. DÉROULEMENT}\end{center}
\par\smallskip\hrule\smallskip
\textbf{Consigne 1}\hfill \textbf{10 min}\par
Développe et réduis $(a+b)^2$, $(a-b)^2$ et $(a+b)(a-b)$, puis formule les égalités obtenues.\par
\textbf{Résultats attendus :}\par Résultat saisi pour la vérification enseignant : les trois identités sont formulées en plénière.\par
\par\smallskip\hrule\smallskip
\textbf{Propriété 2}\hfill \textbf{0 min}\par
Pour tous nombres $a$ et $b$ : $(a+b)^2=a^2+2ab+b^2$ ; $(a-b)^2=a^2-2ab+b^2$ ; $(a+b)(a-b)=a^2-b^2$.\par
\par\smallskip\hrule\smallskip
\textbf{Mise en situation 3 : Situation de départ}\hfill \textbf{0 min}\par
Fofo dispose d'une parcelle rectangulaire dont les côtés sont partagés en longueurs $x$, $z$ d'une part et $y$, $b$ d'autre part. Exprimer son aire de deux façons.\par
\par\smallskip\hrule\smallskip
\textbf{Consigne 4}\hfill \textbf{0 min}\par
Calcule l'aire du grand rectangle, puis la somme des aires des quatre rectangles. Compare les résultats et complète : $(x+z)(y+b)=\ldots$\par
\par\smallskip\hrule\smallskip
\textbf{Trace attendue 5}\hfill \textbf{0 min}\par
$(x+z)(y+b)=xy+xb+zy+zb$. Cette égalité illustre la distributivité de la multiplication sur l'addition.\par
\par\smallskip\hrule\smallskip
\textbf{Application 6}\hfill \textbf{0 min}\par
Développe : $A=(2x+1)^2$, $B=(3x-2)^2$, $C=(4x-3)(4x+3)$. Factorise : $D=x^2+6x+9$, $E=4x^2-9$, $F=4x^2-4x+1$.\par
\par\smallskip\hrule\smallskip
\textbf{Figure 7 : Schéma de la parcelle}\hfill \textbf{5 min}\par
\begin{tikzpicture}[scale=0.75]\draw[thick] (0,0) rectangle (8,4);\draw (5,0)--(5,4);\draw (0,2.5)--(8,2.5);\node at (2.5,1.25) {$xy$};\node at (6.5,1.25) {$zy$};\node at (2.5,3.25) {$xb$};\node at (6.5,3.25) {$zb$};\draw[<->] (0,-0.4)--node[below] {$x$}(5,-0.4);\draw[<->] (5,-0.4)--node[below] {$z$}(8,-0.4);\end{tikzpicture}\par
\end{document}